\documentclass[12pt,a4paper]{article}
\usepackage[T1]{fontenc}
\usepackage[utf8]{inputenc}
\usepackage[polish]{babel}
\usepackage{lmodern}
\usepackage{geometry}
\usepackage{graphicx}
\usepackage{booktabs}
\usepackage{float}
\usepackage{siunitx}
\usepackage{hyperref}
\usepackage{array}
\usepackage{longtable}
\usepackage{caption}
\geometry{margin=2.2cm}
\sisetup{
  output-decimal-marker={,},
  group-separator={\,},
  group-minimum-digits=4
}
\hypersetup{
  colorlinks=true,
  linkcolor=black,
  urlcolor=blue
}

\newcommand{\safeincludegraphics}[2][]{%
  \IfFileExists{#2}{%
    \includegraphics[#1]{#2}%
  }{%
    \fbox{%
      \parbox[c][0.30\textheight][c]{0.92\linewidth}{%
        \centering
        Brak pliku wykresu:\\[0.5em]
        \texttt{#2}%
      }%
    }%
  }%
}

\title{Sprawozdanie z laboratorium 1.}
\author{Jakub Kogut}
\date{\today}

\begin{document}
\maketitle

\section{Cel ćwiczenia}
Celem ćwiczenia było zbadanie działania heurystyki \emph{Local Search} dla euklidesowego
problemu komiwojażera oraz porównanie trzech wariantów:
\begin{enumerate}
    \item pełnego przeszukiwania otoczenia typu \texttt{invert},
    \item wersji przyspieszonej, w której w każdej iteracji rozpatruje się tylko losową próbkę
    sąsiadów,
    \item wariantu dla trzeciego roku, w którym rozwiązanie początkowe budowane jest na bazie
    minimalnego drzewa rozpinającego i przeszukiwania DFS.
\end{enumerate}

Badane instancje pochodzą z kolekcji TSPLIB i odpowiadają pięciu dużym zbiorom:
\texttt{mu1979} (Oman), \texttt{ca4663} (Kanada), \texttt{tz6117} (Tanzania),
\texttt{eg7146} (Egipt) oraz \texttt{ei8246} (Irlandia).

\section{Opis metody}
\subsection*{Reprezentacja rozwiązania}
Rozwiązanie TSP reprezentowane jest jako permutacja wierzchołków. Długość cyklu liczona jest
zgodnie z regułą \texttt{EUC\_2D}, czyli przez zaokrąglenie odległości euklidesowej do najbliższej
liczby całkowitej.

\subsection*{Otoczenie \texttt{invert}}
Podstawowym ruchem lokalnym jest \texttt{invert(i,j)}, czyli odwrócenie fragmentu permutacji
od pozycji \(i\) do pozycji \(j\). W euklidesowym wariancie TSP ruch ten odpowiada klasycznej
operacji 2-opt i bardzo często usuwa skrzyżowania krawędzi.

\subsection*{Warianty eksperymentu}
Dla każdej instancji porównano trzy procedury:
\begin{enumerate}
    \item \textbf{Exercise 1 / full invert}: lokalne przeszukiwanie z pełnym przeglądem
    otoczenia \texttt{invert}.
    \item \textbf{Exercise 2 / sampled invert}: w każdej iteracji oceniana jest jedynie losowa
    próbka sąsiadów, a następnie wybierany jest najlepszy z nich.
    \item \textbf{Exercise 3 / MST seeded}: budowane jest MST, losowany jest wierzchołek startowy,
    tworzony jest cykl z kolejności DFS po drzewie, a następnie uruchamiane jest lokalne
    przeszukiwanie.
\end{enumerate}

\section{Odpowiedzi na pytania z treści zadania}
\subsection*{Dlaczego ruchy typu \texttt{invert} prowadzą do tras bez skrzyżowań?}
W klasycznym euklidesowym TSP przecięcie dwóch krawędzi oznacza, że można zastosować ruch typu
2-opt i zastąpić te krawędzie dwiema innymi, nieskrzyżowanymi. Na mocy nierówności trójkąta taka
zamiana skraca trasę, więc lokalne optimum względem ruchu \texttt{invert} nie powinno zawierać
skrzyżowań.

\subsection*{Dlaczego zaokrąglenia mogą zaburzyć to rozumowanie?}
W używanych danych długości są zaokrąglane do liczb całkowitych. Wtedy może się zdarzyć, że dwie
konfiguracje geometrycznie różniące się długością po zaokrągleniu otrzymają ten sam koszt,
a lokalna poprawa wynikająca z geometrii nie będzie już widoczna w funkcji celu.

\subsection*{Po co używać MST w zadaniu 3?}
Minimalne drzewo rozpinające dostarcza taniej struktury bazowej, która dobrze odwzorowuje
geometrię rozmieszczenia punktów. Cykl otrzymany z DFS po MST zwykle jest wyraźnie lepszym
punktem startowym niż losowa permutacja, więc lokalne przeszukiwanie może szybciej dojść
do rozsądnego rozwiązania.

\section{Wykresy porównawcze}
Poniższe wykresy zbiorcze pokazują trzy najważniejsze miary jakości:
\begin{enumerate}
    \item średnią długość końcowego rozwiązania,
    \item średnią liczbę zaakceptowanych kroków poprawy,
    \item najlepszą uzyskaną długość cyklu.
\end{enumerate}

Osobno pokazano również wagę minimalnego drzewa rozpinającego, ponieważ w zadaniu 3 pełni ono
rolę rozwiązania bazowego i stanowi dobry punkt odniesienia dla jakości inicjalizacji.

\begin{figure}[H]
    \centering
    \safeincludegraphics[width=0.88\textwidth]{plots_lab1/summary/avg_final_length.png}
    \caption{Porównanie średniej długości końcowego rozwiązania dla wszystkich instancji i trzech
    wariantów heurystyki.}
    \label{fig:avg-final}
\end{figure}

\begin{figure}[H]
    \centering
    \safeincludegraphics[width=0.88\textwidth]{plots_lab1/summary/best_length.png}
    \caption{Porównanie najlepszej długości trasy uzyskanej dla każdej instancji.}
    \label{fig:best-length}
\end{figure}

\begin{figure}[H]
    \centering
    \safeincludegraphics[width=0.88\textwidth]{plots_lab1/summary/mst_weight.png}
    \caption{Waga minimalnego drzewa rozpinającego dla badanych instancji.}
    \label{fig:mst-weight}
\end{figure}

\section{Wizualizacje najlepszych tras}
Kolejne rysunki pokazują najlepsze rozwiązania zapisane dla każdej instancji. Każda plansza
zawiera trzy podwykresy odpowiadające trzem wariantom algorytmu, więc można bezpośrednio
porównać wpływ sposobu inicjalizacji i doboru otoczenia na geometrię otrzymanej trasy.

\begin{figure}[H]
    \centering
    \safeincludegraphics[width=\textwidth]{plots_lab1/tours/mu1979_best_tours.png}
    \caption{Najlepsze trasy dla instancji \texttt{mu1979} (Oman).}
\end{figure}

\begin{figure}[H]
    \centering
    \safeincludegraphics[width=\textwidth]{plots_lab1/tours/ca4663_best_tours.png}
    \caption{Najlepsze trasy dla instancji \texttt{ca4663} (Kanada).}
\end{figure}

\begin{figure}[H]
    \centering
    \safeincludegraphics[width=\textwidth]{plots_lab1/tours/tz6117_best_tours.png}
    \caption{Najlepsze trasy dla instancji \texttt{tz6117} (Tanzania).}
\end{figure}

\begin{figure}[H]
    \centering
    \safeincludegraphics[width=\textwidth]{plots_lab1/tours/eg7146_best_tours.png}
    \caption{Najlepsze trasy dla instancji \texttt{eg7146} (Egipt).}
\end{figure}

\begin{figure}[H]
    \centering
    \safeincludegraphics[width=\textwidth]{plots_lab1/tours/ei8246_best_tours.png}
    \caption{Najlepsze trasy dla instancji \texttt{ei8246} (Irlandia).}
\end{figure}

\section{Omówienie wyników}
Najważniejszym wykresem porównawczym jest Rysunek~\ref{fig:avg-final}, ponieważ pokazuje on
rzeczywistą jakość końcową heurystyk po uwzględnieniu przyjętych limitów obliczeniowych.
Jeżeli dla danej instancji wariant MST-seeded osiąga niższą wartość niż oba warianty losowe,
oznacza to, że lepsza inicjalizacja była ważniejsza niż liczba rozpatrywanych sąsiadów.
Jeżeli natomiast pełny wariant \texttt{invert} wygrywa z wersją próbkowaną, to koszt dokładniejszego
przeglądu otoczenia opłacił się jakościowo.

Rysunek~\ref{fig:best-length} pokazuje, który wariant był w stanie znaleźć najlepsze rozwiązanie
w całym eksperymencie dla danej instancji. W praktyce jest to zwykle najważniejsza obserwacja
z punktu widzenia jakości heurystyki. Z kolei wykres wag MST pozwala oszacować, jak trudna jest sama
struktura geometryczna danej instancji i czy rozwiązania z zadania 3 pozostają blisko sensownego
punktu startowego.

Wizualizacje tras są użyteczne z innego powodu niż same wykresy słupkowe. Pozwalają sprawdzić,
czy końcowe rozwiązania zachowują spójną strukturę geometryczną, czy pojawiają się długie
``przeskoki'' przez całą instancję oraz czy wariant MST-seeded rzeczywiście daje lepiej
uporządkowaną trasę już na poziomie kształtu.

\section{Wnioski}
Porównanie trzech wariantów lokalnego przeszukiwania pokazuje dwa praktyczne trade-offy:
\begin{enumerate}
    \item lepsza inicjalizacja może być równie ważna jak dokładniejsze przeszukiwanie otoczenia,
    \item przy bardzo dużych instancjach budżet obliczeniowy silnie wpływa na wynik końcowy,
    więc porównanie algorytmów należy zawsze odnosić do tego samego limitu czasu lub liczby ruchów.
\end{enumerate}

Wariant z pełnym otoczeniem \texttt{invert} stanowi najmocniejszy punkt odniesienia jakościowego,
wariant próbkowany jest kompromisem między czasem a jakością, a wariant MST-seeded pokazuje, że
warto inwestować w lepsze rozwiązanie początkowe. Dalszym naturalnym krokiem byłoby zdjęcie części
ograniczeń lub uczynienie ich parametrami eksperymentu, aby sprawdzić, który z wariantów skaluje się
najlepiej przy rosnącym budżecie obliczeniowym.

\end{document}
